\section{自然試料のδ代表値と基準尺度}

\subsection{平成29年度 問18}
\label{q:h29-reg-applied-18}

\noindent\textbf{原典：} sources/H29/17\_applied\_science.pdf、PDF物理6頁。

\textbf{公式正答：} 3（single）。

\par\textbf{検討未確定：} 自然試料の代表値は地域・時代・生成条件で変わる。全候補の厳密な分布未照合。NBS-127の旧CDT脚注を保持。

\paragraph{問題の要旨} 自然試料のδ値と元素ごとの基準尺度を比較する。

\Needspace{4\baselineskip}\paragraph{確認した条件・記述（独自要約）}\mbox{}\par

\noindent\textbf{判断対象：} 自然試料のδ値と元素ごとの基準尺度を比較する。\par

\noindent\textbf{原典確認：} 全5候補・指数・否定・記号を保存原典22頁画像と照合。\par

\paragraph{問われる核心} 自然試料のδ値と元素ごとの基準尺度を比較する。

\paragraph{必要知識} δ=1000(R試料/R基準−1)で、‰を\%へ誤読しない。

\paragraph{考え方・正答を選ぶ理由（査読待ち初稿）} δ=1000(R試料/R基準−1)で、‰を\%へ誤読しない。海水硫酸の硫黄はCDT/VCDTに対し約+20‰台で0ではなく3。IAEAの海水硫酸由来NBS-127は+20.3‰で旧CDTの脚注もある。これは全海水を同じ厳密数値へ固定する資料ではない。O/C/S/Cl/Nそれぞれの基準を区別する。

\Needspace{4\baselineskip}\paragraph{選択肢の検討（査読待ち初稿）}\mbox{}\par

\noindent\textbf{1：選ばない候補} 海水δ18Oは0‰前後。VSMOW尺度の代表的海水を考える。地域・塩分差はある。\par

\noindent\textbf{2：選ばない候補} 海成炭酸塩δ13Cは0‰前後。VPDB尺度の代表的海成炭酸塩。全地質時代・全生成物の保証ではない。\par

\noindent\textbf{3：選ぶ候補} 海水硫酸δ34Sは0‰前後。海水硫酸由来NBS-127でも約+20.3‰とされ、0ではない。\par

\noindent\textbf{4：選ばない候補} 海水塩化物δ37Clは0‰前後。標準海水塩化物尺度の基準近傍。\par

\noindent\textbf{5：選ばない候補} 大気窒素δ15Nは0‰前後。大気N2を基準とする尺度の零点。\par

\paragraph{再利用する考え方} 軌道と電子状態、格子と基底、単位胞原子数と配位数、定義の分母を区別して候補を全て評価する。

\paragraph{共通理論} \hyperref[theory:isotope:delta]{第\ref*{theory:isotope:delta}節}, \hyperref[theory:isotope:standards]{第\ref*{theory:isotope:standards}節}

\par 参考：\cite{atomic_sulfate,iso_standards1995}。

\par 難易度 2（校準初稿）、優先度未評価。

\section{宇宙線生成核種・露出と堆積年代}

\subsection{平成29年度 問30}
\label{q:h29-reg-applied-30}

\noindent\textbf{原典：} sources/H29/17\_applied\_science.pdf、PDF物理10頁。

\textbf{公式正答：} 3（single）。

\paragraph{問題の要旨} 年代時計のうち、宇宙線で生成される親核種を使う方法を選ぶ。

\Needspace{4\baselineskip}\paragraph{確認した条件・記述（独自要約）}\mbox{}\par

\noindent\textbf{方向：} 宇宙線生成核種を利用する方法として最も適切。\par

\noindent\textbf{対象：} 同位体の標準・分別又は年代。番号・核種・旧標準名称を原典画像で確認。\par

\paragraph{問われる核心} 核種生成と壊変時計・蓄積線量の区別

\paragraph{必要知識} 宇宙線一次粒子と二次中性子、14N(n,p)14C、放射壊変と損傷・捕獲電荷の違い。

\paragraph{考え方・正答を選ぶ理由（査読待ち初稿）} 3の14C法である。生成反応\({}^{14}\mathrm N+n\to{}^{14}\mathrm C+p\)を、形成後の14Cのβ壊変とは区別する。提示中の他法は自然長寿命核種・核分裂飛跡・蓄積線量に基づく。宇宙線が線量に寄与することと、宇宙線生成核種を年代指標にすることを分ける。

\Needspace{4\baselineskip}\paragraph{選択肢の検討（査読待ち初稿）}\mbox{}\par

\noindent\textbf{1：該当しない} U-Pb法。自然界の長寿命U核種の壊変を使い、通常の時計のUを宇宙線生成核種とはしない。\par

\noindent\textbf{2：該当しない} Rb-Sr法。長寿命87Rbの壊変を使い、宇宙線による生成時計ではない。\par

\noindent\textbf{3：該当する} 14C法。大気14Nへ宇宙線の二次中性子が作用して14Cを作る。生物へ入った後の減衰を時計にする。\par

\noindent\textbf{4：該当しない} フィッション・トラック法。自然Uの自発核分裂の飛跡を使う。飛跡一般には宇宙線損傷もあるが、この年代法の時計はU核分裂。\par

\noindent\textbf{5：線量寄与と核種生成が別} 熱ルミネッセンス法。放射線で生じた捕獲電荷の蓄積を発光として測る。宇宙線の線量寄与はあり得るが、宇宙線で作った核種の壊変比を用いる時計ではない。\par

\paragraph{再利用する考え方} 核種の生成、試料への取り込み、移動・分別、閉鎖、壊変を順番に分ける。比の基準と時計の開始条件を特定し、測定値と解釈する年代を区別する。

\paragraph{共通理論} \hyperref[theory:isotope:cosmogenic]{第\ref*{theory:isotope:cosmogenic}節}, \hyperref[theory:isotope:radiocarbon]{第\ref*{theory:isotope:radiocarbon}節}, \hyperref[theory:isotope:method]{第\ref*{theory:isotope:method}節}

\par 参考：\cite{iso_radiometric}。

\par 難易度 2（校準初稿）、優先度未評価。

\section{同位体標準試料}

\subsection{平成23年度 問21}
\label{q:h23-reg-applied-21}

\noindent\textbf{原典：} sources/H23/17\_applied\_science.pdf、PDF物理7頁。

\textbf{公式正答：} 2（single）。

\par\textbf{検討未確定：} 原典のSMOW/PDB/CDTは歴史的対応を問う旧名称である。1995年IAEA勧告は既にVSMOW/VPDB/VCDTと尺度の実現物質を区別しており、出題時に旧隕石・化石標準そのものを現在も測定用一次物質として使うとは説明できない。公式と初回の対応は保存し、尺度と物質・旧名称の省略を査読事項に保持。

\paragraph{問題の要旨} 元素と同位体比の歴史的標準の対応を比較し、炭素と硫黄の基準を取り違えた組合せを選ぶ。

\Needspace{4\baselineskip}\paragraph{確認した条件・記述（独自要約）}\mbox{}\par

\noindent\textbf{方向：} 最も不適切なもの。\par

\noindent\textbf{対象：} 同位体の標準・分別又は年代。番号・核種・旧標準名称を原典画像で確認。\par

\paragraph{問われる核心} 同位体比の基準尺度と標準物質

\paragraph{必要知識} δ=(R/Rstd−1)×1000の定義、H/OのVSMOW、炭酸塩C/OのVPDB、SのVCDT、尺度を実現する標準物質。

\paragraph{考え方・正答を選ぶ理由（査読待ち初稿）} 炭素はPDB/VPDB、硫黄はCDT/VCDT系列という対応なので、CとCDTの2が誤りである。水のHとOはSMOW系列、炭酸塩OはPDB系列を用いる。δ値は比の基準を変えると変わるため、元素だけでなく試料と尺度を明記する。原典旧名称の歴史的対応と、現在の標準物質の実現を分ける。

\Needspace{4\baselineskip}\paragraph{選択肢の検討（査読待ち初稿）}\mbox{}\par

\noindent\textbf{1：対応は意図適切・旧名称注意} HとSMOW。水の2H/1HはSMOW系列（現在のVSMOW-SLAP尺度）。水素をCDT系列としない。\par

\noindent\textbf{2：対応不適切} CとCDT。炭素13C/12CはPDB/VPDB系列である。CDTは硫黄の尺度の起源であり、炭素の標準という対応は誤り。\par

\noindent\textbf{3：対応は意図適切・旧名称注意} OとSMOW。水・ケイ酸塩等の18O/16OにはSMOW/VSMOW系列が使われる。対象により尺度を明記する。\par

\noindent\textbf{4：対応は意図適切・旧名称注意} SとCDT。CDTはCanyon Diablo隕石のトロイライトFeSに由来する硫黄の歴史的基準。実際のVCDT実現は硫化銀IAEA-S-1等を使い、旧CDT物質を無条件に現行標準としない。\par

\noindent\textbf{5：対応は意図適切・旧名称注意} OとPDB。炭酸塩酸素はPDB/VPDB系列でも表す。同じ18Oでも水のVSMOWと炭酸塩のVPDBは数値基準が異なる。\par

\paragraph{再利用する考え方} 核種の生成、試料への取り込み、移動・分別、閉鎖、壊変を順番に分ける。比の基準と時計の開始条件を特定し、測定値と解釈する年代を区別する。

\paragraph{共通理論} \hyperref[theory:isotope:delta]{第\ref*{theory:isotope:delta}節}, \hyperref[theory:isotope:standards]{第\ref*{theory:isotope:standards}節}

\par 参考：\cite{iso_standards1995,iso_vsmow2,iso_vcdt}。

\paragraph{関連問題} \hyperref[q:h30-reg-applied-21]{H30-REG-APPLIED-21}

\par 難易度 2（校準初稿）、優先度未評価。

\subsection{平成30年度 問21}
\label{q:h30-reg-applied-21}

\noindent\textbf{原典：} sources/H30/17\_applied\_science.pdf、PDF物理7頁。

\textbf{公式正答：} 3（single）。

\par\textbf{検討未確定：} 原典のSMOW/PDB/CDTは歴史的対応を問う旧名称である。1995年IAEA勧告は既にVSMOW/VPDB/VCDTと尺度の実現物質を区別しており、出題時に旧隕石・化石標準そのものを現在も測定用一次物質として使うとは説明できない。公式と初回の対応は保存し、尺度と物質・旧名称の省略を査読事項に保持。

\paragraph{問題の要旨} 元素と同位体比の歴史的標準の対応を比較し、炭素と硫黄の基準を取り違えた組合せを選ぶ。

\Needspace{4\baselineskip}\paragraph{確認した条件・記述（独自要約）}\mbox{}\par

\noindent\textbf{方向：} 最も不適切なもの。\par

\noindent\textbf{対象：} 同位体の標準・分別又は年代。番号・核種・旧標準名称を原典画像で確認。\par

\paragraph{問われる核心} 同位体比の基準尺度と標準物質

\paragraph{必要知識} δ=(R/Rstd−1)×1000の定義、H/OのVSMOW、炭酸塩C/OのVPDB、SのVCDT、尺度を実現する標準物質。

\paragraph{考え方・正答を選ぶ理由（査読待ち初稿）} 炭素はPDB/VPDB、硫黄はCDT/VCDT系列という対応なので、CとCDTの3が誤りである。水のHとOはSMOW系列、炭酸塩OはPDB系列を用いる。δ値は比の基準を変えると変わるため、元素だけでなく試料と尺度を明記する。原典旧名称の歴史的対応と、現在の標準物質の実現を分ける。

\Needspace{4\baselineskip}\paragraph{選択肢の検討（査読待ち初稿）}\mbox{}\par

\noindent\textbf{1：対応は意図適切・旧名称注意} SとCDT。CDTはCanyon Diablo隕石のトロイライトFeSに由来する硫黄の歴史的基準。実際のVCDT実現は硫化銀IAEA-S-1等を使い、旧CDT物質を無条件に現行標準としない。\par

\noindent\textbf{2：対応は意図適切・旧名称注意} OとSMOW。水・ケイ酸塩等の18O/16OにはSMOW/VSMOW系列が使われる。対象により尺度を明記する。\par

\noindent\textbf{3：対応不適切} CとCDT。炭素13C/12CはPDB/VPDB系列である。CDTは硫黄の尺度の起源であり、炭素の標準という対応は誤り。\par

\noindent\textbf{4：対応は意図適切・旧名称注意} OとPDB。炭酸塩酸素はPDB/VPDB系列でも表す。同じ18Oでも水のVSMOWと炭酸塩のVPDBは数値基準が異なる。\par

\noindent\textbf{5：対応は意図適切・旧名称注意} HとSMOW。水の2H/1HはSMOW系列（現在のVSMOW-SLAP尺度）。水素をCDT系列としない。\par

\paragraph{再利用する考え方} 核種の生成、試料への取り込み、移動・分別、閉鎖、壊変を順番に分ける。比の基準と時計の開始条件を特定し、測定値と解釈する年代を区別する。

\paragraph{共通理論} \hyperref[theory:isotope:delta]{第\ref*{theory:isotope:delta}節}, \hyperref[theory:isotope:standards]{第\ref*{theory:isotope:standards}節}

\par 参考：\cite{iso_standards1995,iso_vsmow2,iso_vcdt}。

\paragraph{関連問題} \hyperref[q:h23-reg-applied-21]{H23-REG-APPLIED-21}

\par 難易度 2（校準初稿）、優先度未評価。

\section{放射年代測定}

\subsection{平成23年度 問24}
\label{q:h23-reg-applied-24}

\noindent\textbf{原典：} sources/H23/17\_applied\_science.pdf、PDF物理8頁。

\textbf{公式正答：} 4（single）。

\paragraph{問題の要旨} 2万年前の貝化石に対し、試料成分と時計の時間範囲に適合する測定法を選ぶ。

\Needspace{4\baselineskip}\paragraph{確認した条件・記述（独自要約）}\mbox{}\par

\noindent\textbf{試料：} 2万年前の貝化石。海生か淡水生かは原典に指定なし。\par

\noindent\textbf{選択：} K-Ar、Rb-Sr、U-Pb、14C、フィッション・トラックの5法。\par

\noindent\textbf{方向：} 最も適切なもの。\par

\paragraph{問われる核心} 試料と放射性時計の時間尺度の適合

\paragraph{必要知識} 閉鎖系の壊変N=\allowbreak{}N0exp(−λt)、14Cの物理半減期約5730年、生物炭素の交換停止、海洋・淡水リザーバー効果、校準。

\paragraph{考え方・正答を選ぶ理由（査読待ち初稿）} 貝殻を形成した炭素の14Cを用いる4が提示中最も適切である。\(N/N_0=2^{-20000/5730}\simeq0.089\)であり、数万年の測定範囲に合う。他法の対象元素・鉱物と熱履歴はこの若い生物起源炭酸塩の時計に直接対応しない。貝が海生か淡水生か、リザーバー効果・再堆積・炭素交換を確認して暦年を推定する。

\Needspace{4\baselineskip}\paragraph{選択肢の検討（査読待ち初稿）}\mbox{}\par

\noindent\textbf{1：本試料では選ばない} K-Ar法。Kを含む鉱物等の40Kから生成した40Arを測る。2万年前の貝殻炭酸塩の生物由来炭素を直接測る方法としては選ばない。若い火山岩の測定まで絶対不可能という意味ではない。\par

\noindent\textbf{2：本試料では選ばない} Rb-Sr法。Rbを含む鉱物・岩石の87Rb→87Srを用い初生Srと閉鎖条件を検討する。提示の若い貝殻へは14Cほど直接的でない。\par

\noindent\textbf{3：本試料では選ばない} U-Pb法。ジルコン等のUを取り込む鉱物のU→Pb系列に適する。炭酸塩のU系列法一般とU-Pb法を混同せず、提示5法では14Cを選ぶ。\par

\noindent\textbf{4：最も適切（補正は別）} 14C法。生存中取り込まれた14Cが炭素交換停止後に減ることを使う。物理半減期5730年なら2万年後の残存率は約0.089であり測定時間範囲内である。海洋リザーバー・混入・再結晶・暦年校準を別途検討する。\par

\noindent\textbf{5：本試料では選ばない} フィッション・トラック法。Uの自発核分裂による鉱物中の飛跡を用い、加熱で飛跡が消えるため熱履歴に依存する。若い貝殻の炭素を直接測る方法ではない。\par

\paragraph{再利用する考え方} 核種の生成、試料への取り込み、移動・分別、閉鎖、壊変を順番に分ける。比の基準と時計の開始条件を特定し、測定値と解釈する年代を区別する。

\paragraph{共通理論} \hyperref[theory:isotope:decay]{第\ref*{theory:isotope:decay}節}, \hyperref[theory:isotope:radiocarbon]{第\ref*{theory:isotope:radiocarbon}節}, \hyperref[theory:isotope:method]{第\ref*{theory:isotope:method}節}

\par 参考：\cite{iso_radiometric,iso_shell}。

\paragraph{関連問題} \hyperref[q:h26-reg-applied-31]{H26-REG-APPLIED-31}, \hyperref[q:h28-reg-applied-32]{H28-REG-APPLIED-32}

\par 難易度 2（校準初稿）、優先度未評価。

\subsection{平成25年度 問27}
\label{q:h25-reg-applied-27}

\noindent\textbf{原典：} sources/H25/17\_applied\_science.pdf、PDF物理10頁。

\textbf{公式正答：} 2（single）。

\paragraph{問題の要旨} U-Pb年代測定に適する鉱物を、Uの取り込み・初生Pb・閉鎖性から選ぶ。

\Needspace{4\baselineskip}\paragraph{確認した条件・記述（独自要約）}\mbox{}\par

\noindent\textbf{方向：} U-Pbに最も適した鉱物。\par

\noindent\textbf{対象：} 同位体の標準・分別又は年代。番号・核種・旧標準名称を原典画像で確認。\par

\paragraph{問われる核心} 放射性時計を保持する鉱物の選択

\paragraph{必要知識} ジルコンZrSiO4、UとZrの置換、放射起源Pbと初生共通Pb、閉鎖系、U-Pb二系列と継承核。

\paragraph{考え方・正答を選ぶ理由（査読待ち初稿）} 2のジルコンを選ぶ。Uを取り込み、初生Pbが比較的少なく、結晶が地質過程を通じてU-Pb系を保ちやすいためである。238U→206Pbと235U→207Pbから独立な年代を比べ、同じ年代かどうかで閉鎖系の仮定を点検する。古い継承核を持つ粒子は岩体の定置年代と一致しない場合がある。

\Needspace{4\baselineskip}\paragraph{選択肢の検討（査読待ち初稿）}\mbox{}\par

\noindent\textbf{1：代表試料ではない} 黒雲母。Kを含む雲母でK-ArやRb-Srの対象になり得るが、提示中のU-Pbの代表試料ではない。\par

\noindent\textbf{2：最適} ジルコン。ZrSiO4中でUがZr位置へ取り込まれ、通常初生Pbは少ない。抵抗性が高くU-Pb年代に適する。継承核・Pb損失・共通Pbは別途検査する。\par

\noindent\textbf{3：代表試料ではない} 白雲母。Kを含む雲母でK-ArやRb-Sr等に利用するが、U-Pbの代表的選択肢はジルコンである。\par

\noindent\textbf{4：代表試料ではない} 海緑石。Kを含む堆積性鉱物でK-Ar等の候補になり得る。生成・変質条件を要し、提示中のU-Pb最適鉱物ではない。\par

\noindent\textbf{5：代表試料ではない} カリ長石。Kに富む長石で初生Pbを含み得る。Uを集めるジルコンのようなU-Pbの代表鉱物ではない。\par

\paragraph{再利用する考え方} 核種の生成、試料への取り込み、移動・分別、閉鎖、壊変を順番に分ける。比の基準と時計の開始条件を特定し、測定値と解釈する年代を区別する。

\paragraph{共通理論} \hyperref[theory:isotope:upb]{第\ref*{theory:isotope:upb}節}, \hyperref[theory:isotope:closure]{第\ref*{theory:isotope:closure}節}, \hyperref[theory:isotope:method]{第\ref*{theory:isotope:method}節}

\par 参考：\cite{iso_zircon,iso_radiometric}。

\par 難易度 2（校準初稿）、優先度未評価。

\subsection{平成26年度 問31}
\label{q:h26-reg-applied-31}

\noindent\textbf{原典：} sources/H26/17\_applied\_science.pdf、PDF物理13頁。

\textbf{公式正答：} 4（single）。

\paragraph{問題の要旨} 2万年前の貝化石に対し、試料成分と時計の時間範囲に適合する測定法を選ぶ。

\Needspace{4\baselineskip}\paragraph{確認した条件・記述（独自要約）}\mbox{}\par

\noindent\textbf{試料：} 2万年前の貝化石。海生か淡水生かは原典に指定なし。\par

\noindent\textbf{選択：} K-Ar、Rb-Sr、U-Pb、14C、フィッション・トラックの5法。\par

\noindent\textbf{方向：} 最も適切なもの。\par

\paragraph{問われる核心} 試料と放射性時計の時間尺度の適合

\paragraph{必要知識} 閉鎖系の壊変N=\allowbreak{}N0exp(−λt)、14Cの物理半減期約5730年、生物炭素の交換停止、海洋・淡水リザーバー効果、校準。

\paragraph{考え方・正答を選ぶ理由（査読待ち初稿）} 貝殻を形成した炭素の14Cを用いる4が提示中最も適切である。\(N/N_0=2^{-20000/5730}\simeq0.089\)であり、数万年の測定範囲に合う。他法の対象元素・鉱物と熱履歴はこの若い生物起源炭酸塩の時計に直接対応しない。貝が海生か淡水生か、リザーバー効果・再堆積・炭素交換を確認して暦年を推定する。

\Needspace{4\baselineskip}\paragraph{選択肢の検討（査読待ち初稿）}\mbox{}\par

\noindent\textbf{1：本試料では選ばない} K-Ar法。Kを含む鉱物等の40Kから生成した40Arを測る。2万年前の貝殻炭酸塩の生物由来炭素を直接測る方法としては選ばない。若い火山岩の測定まで絶対不可能という意味ではない。\par

\noindent\textbf{2：本試料では選ばない} Rb-Sr法。Rbを含む鉱物・岩石の87Rb→87Srを用い初生Srと閉鎖条件を検討する。提示の若い貝殻へは14Cほど直接的でない。\par

\noindent\textbf{3：本試料では選ばない} U-Pb法。ジルコン等のUを取り込む鉱物のU→Pb系列に適する。炭酸塩のU系列法一般とU-Pb法を混同せず、提示5法では14Cを選ぶ。\par

\noindent\textbf{4：最も適切（補正は別）} 14C法。生存中取り込まれた14Cが炭素交換停止後に減ることを使う。物理半減期5730年なら2万年後の残存率は約0.089であり測定時間範囲内である。海洋リザーバー・混入・再結晶・暦年校準を別途検討する。\par

\noindent\textbf{5：本試料では選ばない} フィッション・トラック法。Uの自発核分裂による鉱物中の飛跡を用い、加熱で飛跡が消えるため熱履歴に依存する。若い貝殻の炭素を直接測る方法ではない。\par

\paragraph{再利用する考え方} 核種の生成、試料への取り込み、移動・分別、閉鎖、壊変を順番に分ける。比の基準と時計の開始条件を特定し、測定値と解釈する年代を区別する。

\paragraph{共通理論} \hyperref[theory:isotope:decay]{第\ref*{theory:isotope:decay}節}, \hyperref[theory:isotope:radiocarbon]{第\ref*{theory:isotope:radiocarbon}節}, \hyperref[theory:isotope:method]{第\ref*{theory:isotope:method}節}

\par 参考：\cite{iso_radiometric,iso_shell}。

\paragraph{関連問題} \hyperref[q:h23-reg-applied-24]{H23-REG-APPLIED-24}

\par 難易度 2（校準初稿）、優先度未評価。

\subsection{平成28年度 問32}
\label{q:h28-reg-applied-32}

\noindent\textbf{原典：} sources/H28/17\_applied\_science.pdf、PDF物理13頁。

\textbf{公式正答：} 1（single）。

\paragraph{問題の要旨} 白亜紀試料に対し、放射性時計の時間範囲が短すぎる方法を選ぶ。

\Needspace{4\baselineskip}\paragraph{確認した条件・記述（独自要約）}\mbox{}\par

\noindent\textbf{対象：} 白亜紀の地質試料。具体的鉱物は指定されていない。\par

\noindent\textbf{方向：} 測定方法として最も不適切。\par

\paragraph{問われる核心} 年代範囲と時計の感度

\paragraph{必要知識} 壊変率λ=ln2/T、検出限界、14Cの数万年範囲、K-Ar/U-Pb/Rb-Srと飛跡法の試料・閉鎖条件。

\paragraph{考え方・正答を選ぶ理由（査読待ち初稿）} 14Cの1が不適切である。数万年程度の時計を数千万〜億年の試料へ使うと、初生14Cは事実上失われる。他の4法は適切な鉱物・岩石と閉鎖条件を選べばこの時間尺度の年代又は冷却履歴を扱える。長半減期だけで万能とせず、何の出来事から時計が閉じたかを確かめる。

\Needspace{4\baselineskip}\paragraph{選択肢の検討（査読待ち初稿）}\mbox{}\par

\noindent\textbf{1：不適切} 14C法。白亜紀の数千万〜億年の時間は、約5730年という半減期を何千回も超える。測定可能な14Cは背景・汚染に支配され、直接年代時計として不適切。\par

\noindent\textbf{2：適切（条件付き）} K-Ar法。40Kの長い半減期と蓄積した放射起源40Arを使う。K含有鉱物を含む適切な岩石を選び、Ar損失や過剰Arを検査する。\par

\noindent\textbf{3：適切（条件付き）} U-Pb法。白亜紀ジルコン等でU→Pbの二系列を使える。鉱物選択と初生Pb、Pb損失、継承核を検査する。\par

\noindent\textbf{4：適切（条件付き）} Rb-Sr法。87Rb→87Srの比と複数の同源試料により初生Srを分ける。白亜紀の時間自体は適用候補だが、あらゆる試料に使えるとはしない。\par

\noindent\textbf{5：適切（条件付き）} フィッション・トラック法。アパタイト・ジルコン等で238U自発核分裂の損傷を測る。白亜紀の冷却・熱履歴を測れるが、形成年代と必ず一致するとはしない。\par

\paragraph{再利用する考え方} 核種の生成、試料への取り込み、移動・分別、閉鎖、壊変を順番に分ける。比の基準と時計の開始条件を特定し、測定値と解釈する年代を区別する。

\paragraph{共通理論} \hyperref[theory:isotope:decay]{第\ref*{theory:isotope:decay}節}, \hyperref[theory:isotope:method]{第\ref*{theory:isotope:method}節}, \hyperref[theory:isotope:closure]{第\ref*{theory:isotope:closure}節}

\par 参考：\cite{iso_radiometric,iso_zircon}。

\paragraph{関連問題} \hyperref[q:h23-reg-applied-24]{H23-REG-APPLIED-24}, \hyperref[q:r01-ret-applied-35]{R01-RET-APPLIED-35}

\par 難易度 2（校準初稿）、優先度未評価。

\subsection{令和元年度再試験 問35}
\label{q:r01-ret-applied-35}

\noindent\textbf{原典：} sources/R01\_retest/17\_applied\_science.pdf、PDF物理11頁。

\textbf{公式正答：} 5（single）。

\paragraph{問題の要旨} 白亜紀試料に対し、放射性時計の時間範囲が短すぎる方法を選ぶ。

\Needspace{4\baselineskip}\paragraph{確認した条件・記述（独自要約）}\mbox{}\par

\noindent\textbf{対象：} 白亜紀の地質試料。具体的鉱物は指定されていない。\par

\noindent\textbf{方向：} 測定方法として最も不適切。\par

\paragraph{問われる核心} 年代範囲と時計の感度

\paragraph{必要知識} 壊変率λ=ln2/T、検出限界、14Cの数万年範囲、K-Ar/U-Pb/Rb-Srと飛跡法の試料・閉鎖条件。

\paragraph{考え方・正答を選ぶ理由（査読待ち初稿）} 14Cの5が不適切である。数万年程度の時計を数千万〜億年の試料へ使うと、初生14Cは事実上失われる。他の4法は適切な鉱物・岩石と閉鎖条件を選べばこの時間尺度の年代又は冷却履歴を扱える。長半減期だけで万能とせず、何の出来事から時計が閉じたかを確かめる。

\Needspace{4\baselineskip}\paragraph{選択肢の検討（査読待ち初稿）}\mbox{}\par

\noindent\textbf{1：適切（条件付き）} K-Ar法。40Kの長い半減期と蓄積した放射起源40Arを使う。K含有鉱物を含む適切な岩石を選び、Ar損失や過剰Arを検査する。\par

\noindent\textbf{2：適切（条件付き）} U-Pb法。白亜紀ジルコン等でU→Pbの二系列を使える。鉱物選択と初生Pb、Pb損失、継承核を検査する。\par

\noindent\textbf{3：適切（条件付き）} Rb-Sr法。87Rb→87Srの比と複数の同源試料により初生Srを分ける。白亜紀の時間自体は適用候補だが、あらゆる試料に使えるとはしない。\par

\noindent\textbf{4：適切（条件付き）} フィッション・トラック法。アパタイト・ジルコン等で238U自発核分裂の損傷を測る。白亜紀の冷却・熱履歴を測れるが、形成年代と必ず一致するとはしない。\par

\noindent\textbf{5：不適切} 14C法。白亜紀の数千万〜億年の時間は、約5730年という半減期を何千回も超える。測定可能な14Cは背景・汚染に支配され、直接年代時計として不適切。\par

\paragraph{再利用する考え方} 核種の生成、試料への取り込み、移動・分別、閉鎖、壊変を順番に分ける。比の基準と時計の開始条件を特定し、測定値と解釈する年代を区別する。

\paragraph{共通理論} \hyperref[theory:isotope:decay]{第\ref*{theory:isotope:decay}節}, \hyperref[theory:isotope:method]{第\ref*{theory:isotope:method}節}, \hyperref[theory:isotope:closure]{第\ref*{theory:isotope:closure}節}

\par 参考：\cite{iso_radiometric,iso_zircon}。

\paragraph{関連問題} \hyperref[q:h28-reg-applied-32]{H28-REG-APPLIED-32}, \hyperref[q:r04-reg-applied-25]{R04-REG-APPLIED-25}

\par 難易度 2（校準初稿）、優先度未評価。

\subsection{令和4年度 問25}
\label{q:r04-reg-applied-25}

\noindent\textbf{原典：} sources/R04/17\_applied\_science.pdf、PDF物理9頁。

\textbf{公式正答：} 3（single）。

\paragraph{問題の要旨} 白亜紀試料に対し、放射性時計の時間範囲が短すぎる方法を選ぶ。

\Needspace{4\baselineskip}\paragraph{確認した条件・記述（独自要約）}\mbox{}\par

\noindent\textbf{対象：} 白亜紀の地質試料。具体的鉱物は指定されていない。\par

\noindent\textbf{方向：} 測定方法として最も不適切。\par

\paragraph{問われる核心} 年代範囲と時計の感度

\paragraph{必要知識} 壊変率λ=ln2/T、検出限界、14Cの数万年範囲、K-Ar/U-Pb/Rb-Srと飛跡法の試料・閉鎖条件。

\paragraph{考え方・正答を選ぶ理由（査読待ち初稿）} 14Cの3が不適切である。数万年程度の時計を数千万〜億年の試料へ使うと、初生14Cは事実上失われる。他の4法は適切な鉱物・岩石と閉鎖条件を選べばこの時間尺度の年代又は冷却履歴を扱える。長半減期だけで万能とせず、何の出来事から時計が閉じたかを確かめる。

\Needspace{4\baselineskip}\paragraph{選択肢の検討（査読待ち初稿）}\mbox{}\par

\noindent\textbf{1：適切（条件付き）} U-Pb法。白亜紀ジルコン等でU→Pbの二系列を使える。鉱物選択と初生Pb、Pb損失、継承核を検査する。\par

\noindent\textbf{2：適切（条件付き）} K-Ar法。40Kの長い半減期と蓄積した放射起源40Arを使う。K含有鉱物を含む適切な岩石を選び、Ar損失や過剰Arを検査する。\par

\noindent\textbf{3：不適切} 14C法。白亜紀の数千万〜億年の時間は、約5730年という半減期を何千回も超える。測定可能な14Cは背景・汚染に支配され、直接年代時計として不適切。\par

\noindent\textbf{4：適切（条件付き）} Rb-Sr法。87Rb→87Srの比と複数の同源試料により初生Srを分ける。白亜紀の時間自体は適用候補だが、あらゆる試料に使えるとはしない。\par

\noindent\textbf{5：適切（条件付き）} フィッション・トラック法。アパタイト・ジルコン等で238U自発核分裂の損傷を測る。白亜紀の冷却・熱履歴を測れるが、形成年代と必ず一致するとはしない。\par

\paragraph{再利用する考え方} 核種の生成、試料への取り込み、移動・分別、閉鎖、壊変を順番に分ける。比の基準と時計の開始条件を特定し、測定値と解釈する年代を区別する。

\paragraph{共通理論} \hyperref[theory:isotope:decay]{第\ref*{theory:isotope:decay}節}, \hyperref[theory:isotope:method]{第\ref*{theory:isotope:method}節}, \hyperref[theory:isotope:closure]{第\ref*{theory:isotope:closure}節}

\par 参考：\cite{iso_radiometric,iso_zircon}。

\paragraph{関連問題} \hyperref[q:r01-ret-applied-35]{R01-RET-APPLIED-35}

\par 難易度 2（校準初稿）、優先度未評価。

\section{水の放射性トレーサーと滞留時間}

\subsection{令和7年度 問22}
\label{q:r07-reg-applied-22}

\noindent\textbf{原典：} sources/R07/17\_applied\_science.pdf、PDF物理8頁。

\textbf{公式正答：} 4（single）。

\par\textbf{検討未確定：} 初回・公式4を保持。原典はカルサイトに含まれるClへの宇宙線照射による36Cl生成と書く。石灰岩主要経路はCa核破砕だが、微量35Clの二次中性子捕獲経路も存在する。ClをCaへ原典校正せず、Cl経路を不可能とも説明しない。一意な科学的誤り肢かは査読待ち。

\paragraph{問題の要旨} 水時計と宇宙線生成核種の標的元素・輸送を区別して各測定法の前提を検討する。

\Needspace{4\baselineskip}\paragraph{確認した条件・記述（独自要約）}\mbox{}\par

\noindent\textbf{方向：} 最も不適切なもの。\par

\noindent\textbf{対象：} 同位体の標準・分別又は年代。番号・核種・旧標準名称を原典画像で確認。\par

\paragraph{問われる核心} 自然水の親娘時計と宇宙線標的反応

\paragraph{必要知識} 親娘3H/3Heの質量収支、炭素反応補正、10Be大気生成と沈着、Ca核破砕と35Cl中性子捕獲、UThα壊変由来4Heの供給。

\paragraph{考え方・正答を選ぶ理由（査読待ち初稿）} 公式4と初回4は一致する。石灰岩表面の36ClではCa核破砕が主要経路であり、原典がClをカルサイトの標的として挙げる点が意図された誤りと考えられる。一方、35Cl(n,γ)36Clという二次中性子捕獲は実在する。微量Clを含む試料まで含めた原典の広い表現に科学的一意性を与えず、公式と疑義を併記する。他の方法も閉鎖・混合・生成入力の補正を条件とする。

\Needspace{4\baselineskip}\paragraph{選択肢の検討（査読待ち初稿）}\mbox{}\par

\noindent\textbf{1：適切（補正・閉鎖条件）} 3Hのβ壊変で生じる3Heの蓄積と3H/3Heから水塊年代。。大気との交換停止後、3Hから生じる放射起源3Heを分けて親娘比から見かけ年代を求める。初生大気He、過剰空気、地殻・マントルHe、混合とガス損失を補正する。\par

\noindent\textbf{2：適切（反応補正）} 古い地下水14Cは地化学反応の無機炭素値変化を補正。。炭酸塩溶解等の古い炭素が14Cを希釈し、反応・交換・混合で年代指標が変わるため補正が必要。測る炭素の見かけ年代と水の移動時間は同一とは限らない。\par

\noindent\textbf{3：適切（入力・閉鎖条件）} 大気N/Oの核破砕で作る10Beがエアロゾルに付着し堆積速度へ利用。。大気宇宙線核破砕で作った10Beが粒子に付着し湿性・乾性沈着で堆積物へ入る。初期入力、混合・再移動と閉鎖性を検討して堆積履歴を推定する。\par

\noindent\textbf{4：意図不適切・生成経路疑義} カルサイト中Clへ宇宙線が照射して36Clを作り石灰岩削剥を推定。。石灰岩ではカルサイトの主成分Caの核破砕が主要な36Cl生成経路であり、原典Cl標的は意図上取り違えと考えられる。ただし微量35Clへの二次中性子捕獲でも36Clを作る実在経路がある。原典を「Ca」に訂正せず、Cl経路を不可能とも断定しない。\par

\noindent\textbf{5：適切（供給モデル条件）} 岩石U/Thのα壊変で作る4Heが地下水に蓄積し滞留推定。。α粒子は4He核で、U/Th系列由来Heが溶存水へ供給される。生成供給速度・地殻深部流束・拡散漏れ・混合等のモデルを要し、濃度だけから常に一意な年代を得るわけではない。\par

\paragraph{再利用する考え方} 核種の生成、試料への取り込み、移動・分別、閉鎖、壊変を順番に分ける。比の基準と時計の開始条件を特定し、測定値と解釈する年代を区別する。

\paragraph{共通理論} \hyperref[theory:isotope:tritium]{第\ref*{theory:isotope:tritium}節}, \hyperref[theory:isotope:deadcarbon]{第\ref*{theory:isotope:deadcarbon}節}, \hyperref[theory:isotope:helium]{第\ref*{theory:isotope:helium}節}, \hyperref[theory:isotope:cosmogenic]{第\ref*{theory:isotope:cosmogenic}節}, \hyperref[theory:isotope:cl36]{第\ref*{theory:isotope:cl36}節}

\par 参考：\cite{iso_tritium_clock,iso_groundwater,iso_deadcarbon,iso_be10,iso_cl36,iso_cl36_uw}。

\paragraph{関連問題} \hyperref[q:h27-reg-applied-25]{H27-REG-APPLIED-25}

\par 難易度 3（校準初稿）、優先度未評価。

\section{自然水の同位体分別・気候と水岩反応}

\subsection{平成24年度 問32}
\label{q:h24-reg-applied-32}

\noindent\textbf{原典：} sources/H24/17\_applied\_science.pdf、PDF物理9頁。

\textbf{公式正答：} 3（single）。

\paragraph{問題の要旨} 自然水・地下水炭素・炭酸塩の同位体尺度と、蒸発・氷期の分別を比較する。

\Needspace{4\baselineskip}\paragraph{確認した条件・記述（独自要約）}\mbox{}\par

\noindent\textbf{方向：} 最も不適切なもの。\par

\noindent\textbf{対象：} 同位体の標準・分別又は年代。番号・核種・旧標準名称を原典画像で確認。\par

\paragraph{問われる核心} 同位体の基準・混合と気候分別

\paragraph{必要知識} δの基準、水と炭酸塩の尺度、残留水の重同位体濃縮、氷床への軽同位体隔離、炭素混合と見かけ14C年代。

\paragraph{考え方・正答を選ぶ理由（査読待ち初稿）} 3の「SMOWのみ」が不適切である。水・ケイ酸塩と炭酸塩では慣用尺度が異なり、炭酸塩酸素はPDB/VPDB系列を使う。他肢は蒸発残留、古い炭素による希釈、氷床と炭酸塩分別の作用を、それぞれ適用条件を付して理解する。海水が完全一様、全熱水が負値、全有機物が必ず老齢化という全称命題にはしない。

\Needspace{4\baselineskip}\paragraph{選択肢の検討（査読待ち初稿）}\mbox{}\par

\noindent\textbf{1：適切（条件付き）} 海水δ18Oの地域差は天水・熱水より小さく、後者は負値が多い。。外洋海水は天水に比べ同位体組成の変動が小さいという一般的比較である。蒸発・降水・氷融解・河川流入で地域差はある。熱水も天水起源なら負値を取り得るが、水岩交換・マグマ性流体で変わるため全海水・全熱水の固定値とはしない。\par

\noindent\textbf{2：適切（条件付き）} 有機物分解の炭素添加で地下水14C見かけ年代が古くなる場合。。14Cに乏しい古い有機物又は炭酸塩から炭素が加わると、炭素量に対する14Cが薄まり未補正見かけ年代は古くなる。「ことがある」という原典の条件付き主張である。若い有機炭素なら若返る場合もあり、全有機物に一般化しない。\par

\noindent\textbf{3：不適切} 水・珪酸塩岩・炭酸塩岩は全て酸素SMOWのみ。。炭酸塩の酸素比はPDB/VPDB系列でも表示する。全対象にSMOWのみという限定は誤りである。異なる尺度のδ数値はそのまま比較しない。\par

\noindent\textbf{4：適切（条件付き）} 蒸発残留水にDと18Oが濃縮。。軽い水分子が蒸発しやすく、残った水はD（2H）・18Oに富む傾向となる。供給水や湿度・同位体交換による例外を分ける。蒸気が重くなるという説明へ逆転しない。\par

\noindent\textbf{5：適切（条件付き）} 有孔虫δ18Oは氷期の方が間氷期より大きい。。氷期は16Oに富む水が陸上氷に隔離され、海水が18Oに富む。低温ほど炭酸塩と水の分別も殻18Oを増やす方向であり、有孔虫δ18Oは通常氷期に大きい。局所塩分・水温・生物効果・続成変質の補正は別に要する。\par

\paragraph{再利用する考え方} 核種の生成、試料への取り込み、移動・分別、閉鎖、壊変を順番に分ける。比の基準と時計の開始条件を特定し、測定値と解釈する年代を区別する。

\paragraph{共通理論} \hyperref[theory:isotope:delta]{第\ref*{theory:isotope:delta}節}, \hyperref[theory:isotope:standards]{第\ref*{theory:isotope:standards}節}, \hyperref[theory:isotope:fractionation]{第\ref*{theory:isotope:fractionation}節}, \hyperref[theory:isotope:climate]{第\ref*{theory:isotope:climate}節}, \hyperref[theory:isotope:deadcarbon]{第\ref*{theory:isotope:deadcarbon}節}

\par 参考：\cite{iso_standards1995,iso_oxygen,iso_deadcarbon,iso_hydrothermal}。

\paragraph{関連問題} \hyperref[q:r02-reg-applied-20]{R02-REG-APPLIED-20}

\par 難易度 3（校準初稿）、優先度未評価。

\subsection{平成27年度 問25}
\label{q:h27-reg-applied-25}

\noindent\textbf{原典：} sources/H27/17\_applied\_science.pdf、PDF物理10頁。

\textbf{公式正答：} 3（single）。

\paragraph{問題の要旨} 水の同位体比の基準と天水線の軸、核実験トリチウムと分別を比較する。

\Needspace{4\baselineskip}\paragraph{確認した条件・記述（独自要約）}\mbox{}\par

\noindent\textbf{方向：} 最も不適切なもの。\par

\noindent\textbf{対象：} 同位体の標準・分別又は年代。番号・核種・旧標準名称を原典画像で確認。\par

\paragraph{問われる核心} 安定同位体の起源指標と放射性水時計

\paragraph{必要知識} δ比の定義、全球・地域天水線、D=\allowbreak{}2H、核実験3H入力、壊変と混合。

\paragraph{考え方・正答を選ぶ理由（査読待ち初稿）} 3が軸を取り違えている。天水線は\(\delta^2\mathrm H\simeq8\delta^{18}\mathrm O+10\)で、水素と酸素を結ぶ。δは千分率偏差であり濃度ではない。3Hの入力パルス・壊変は年代、安定H/O比は起源・分別を主に追い、両者を区別する。

\Needspace{4\baselineskip}\paragraph{選択肢の検討（査読待ち初稿）}\mbox{}\par

\noindent\textbf{1：適切（条件付き）} 各種同位体から土壌水・地下水・流出水の滞留を推定。。放射性同位体の壊変、安定同位体の入力変動と混合モデルを組み合わせ、適切な条件で滞留を推定できる。安定比だけが単調な年齢時計になるわけではない。\par

\noindent\textbf{2：適切（条件付き）} H/Oのδは標準比からの千分偏差、標準平均海水等。。δ=(R/Rstd−1)×1000である。原典旧SMOW名称をVSMOW-SLAP尺度とその実現物質へ接続する。\par

\noindent\textbf{3：不適切} δ18O対δ17Oを天水線と呼ぶ。。天水線は水素δD（δ2H）と酸素δ18Oの関係である。酸素17/18の相関には別の意味があり、それを天水線と呼ばない。\par

\noindent\textbf{4：適切（条件付き）} 1952年以降核実験トリチウム増加から地下水涵養を推定。。核実験が作った1950年代以降の入力履歴を使う。減衰と混合を補正し、同じトリチウム濃度に複数涵養年が対応し得る。\par

\noindent\textbf{5：適切（条件付き）} 蒸発凝縮の分別でδD・δ18Oが変わる。。質量差による平衡・速度論的分別が蒸発凝縮の履歴を残す。混合・蒸発線と地域天水線を分ける。\par

\paragraph{再利用する考え方} 核種の生成、試料への取り込み、移動・分別、閉鎖、壊変を順番に分ける。比の基準と時計の開始条件を特定し、測定値と解釈する年代を区別する。

\paragraph{共通理論} \hyperref[theory:isotope:delta]{第\ref*{theory:isotope:delta}節}, \hyperref[theory:isotope:meteoric]{第\ref*{theory:isotope:meteoric}節}, \hyperref[theory:isotope:tritium]{第\ref*{theory:isotope:tritium}節}

\par 参考：\cite{gc_meteoric,iso_tritium_clock,iso_vsmow2}。

\paragraph{関連問題} \hyperref[q:r05-reg-applied-21]{R05-REG-APPLIED-21}, \hyperref[q:r07-reg-applied-22]{R07-REG-APPLIED-22}

\par 難易度 3（校準初稿）、優先度未評価。

\subsection{令和2年度 問20}
\label{q:r02-reg-applied-20}

\noindent\textbf{原典：} sources/R02/17\_applied\_science.pdf、PDF物理7頁。

\textbf{公式正答：} 3（single）。

\par\textbf{検討未確定：} 公式と初回は一致。原典の硫黄「国際標準物質CDTを用いて表記」は、出題前1995年のVCDT/IAEA-S-1勧告と物質・尺度の区別を省く。旧名称を現在の測定用標準物質へそのまま昇格しない。氷期δ大小の誤りは独立に説明できるが、硫黄肢の文字通りの妥当性は査読事項。

\paragraph{問題の要旨} 氷期と間氷期の有孔虫酸素比の大小を、蒸発・炭素混合・歴史的硫黄尺度から区別する。

\Needspace{4\baselineskip}\paragraph{確認した条件・記述（独自要約）}\mbox{}\par

\noindent\textbf{方向：} 最も不適切なもの。\par

\noindent\textbf{対象：} 同位体の標準・分別又は年代。番号・核種・旧標準名称を原典画像で確認。\par

\paragraph{問われる核心} 氷床・炭酸塩分別と同位体尺度の時点

\paragraph{必要知識} 海洋と氷の同位体収支、温度依存の炭酸塩分別、古い炭素混合、CDTとVCDT・IAEA-S-1の違い。

\paragraph{考え方・正答を選ぶ理由（査読待ち初稿）} 氷期より間氷期のδ18Oが大きいとする3が意図された不適切肢である。氷床へ軽い水を隔離する作用と、低温で殻が重い酸素を取り込みやすい分別が、通常ともに氷期の殻δ18Oを増す。他肢の条件を分け、硫黄CDTは歴史的起源として説明しVCDT尺度の実現を併記する。

\Needspace{4\baselineskip}\paragraph{選択肢の検討（査読待ち初稿）}\mbox{}\par

\noindent\textbf{1：適切（条件付き）} 蒸発残留水はDと18Oに富む。。軽い水分子が蒸発しやすく、残った水はD（2H）・18Oに富む傾向となる。供給水や湿度・同位体交換による例外を分ける。蒸気が重くなるという説明へ逆転しない。\par

\noindent\textbf{2：適切（条件付き）} 地下水へ有機物分解炭素が加わると14C見かけ年代が古くなる場合。。14Cに乏しい古い有機物又は炭酸塩から炭素が加わると、炭素量に対する14Cが薄まり未補正見かけ年代は古くなる。「ことがある」という原典の条件付き主張である。若い有機炭素なら若返る場合もあり、全有機物に一般化しない。\par

\noindent\textbf{3：意図不適切（大小逆）} 有孔虫δ18Oは氷期より間氷期に大きい。。氷期は16Oに富む水が陸上氷に隔離され、海水が18Oに富む。低温ほど炭酸塩と水の分別も殻18Oを増やす方向であり、有孔虫δ18Oは通常氷期に大きい。局所塩分・水温・生物効果・続成変質の補正は別に要する。原典はこの大小を逆にしている。\par

\noindent\textbf{4：旧尺度として意図適切・名称疑義} 硫黄標準はCDTトロイライト。。CDTはCanyon Diablo隕石のトロイライトFeSに由来する硫黄の歴史的基準。実際のVCDT実現は硫化銀IAEA-S-1等を使い、旧CDT物質を無条件に現行標準としない。\par

\noindent\textbf{5：適切（条件付き）} 海水δ18Oの変動は天水・熱水より小さく後者は負値が多い。。外洋海水は天水に比べ同位体組成の変動が小さいという一般的比較である。蒸発・降水・氷融解・河川流入で地域差はある。熱水も天水起源なら負値を取り得るが、水岩交換・マグマ性流体で変わるため全海水・全熱水の固定値とはしない。\par

\paragraph{再利用する考え方} 核種の生成、試料への取り込み、移動・分別、閉鎖、壊変を順番に分ける。比の基準と時計の開始条件を特定し、測定値と解釈する年代を区別する。

\paragraph{共通理論} \hyperref[theory:isotope:climate]{第\ref*{theory:isotope:climate}節}, \hyperref[theory:isotope:fractionation]{第\ref*{theory:isotope:fractionation}節}, \hyperref[theory:isotope:deadcarbon]{第\ref*{theory:isotope:deadcarbon}節}, \hyperref[theory:isotope:standards]{第\ref*{theory:isotope:standards}節}

\par 参考：\cite{iso_oxygen,iso_deadcarbon,iso_standards1995,iso_vcdt,iso_hydrothermal}。

\paragraph{関連問題} \hyperref[q:h24-reg-applied-32]{H24-REG-APPLIED-32}, \hyperref[q:r06-reg-applied-21]{R06-REG-APPLIED-21}

\par 難易度 3（校準初稿）、優先度未評価。

\subsection{令和5年度 問21}
\label{q:r05-reg-applied-21}

\noindent\textbf{原典：} sources/R05/17\_applied\_science.pdf、PDF物理7頁。

\textbf{公式正答：} 1（single）。

\paragraph{問題の要旨} 安定な重水素と放射性トリチウムを区別し、水の分別・天水線・熱水酸素シフトを読む。

\Needspace{4\baselineskip}\paragraph{確認した条件・記述（独自要約）}\mbox{}\par

\noindent\textbf{方向：} 最も不適切なもの。\par

\noindent\textbf{対象：} 同位体の標準・分別又は年代。番号・核種・旧標準名称を原典画像で確認。\par

\paragraph{問われる核心} 核種の安定性と水の同位体分別

\paragraph{必要知識} 2Hと3Hの核構造・安定性、δ定義、天水線、蒸発凝縮・水岩同位体交換。

\paragraph{考え方・正答を選ぶ理由（査読待ち初稿）} 1が2Hと3Hを混同している。D=\allowbreak{}2Hは安定、T=\allowbreak{}3Hは半減期約12.3年で3Heへβ壊変する。安定H/Oの比は蒸発凝縮や水岩交換の履歴を追う指標であり、この壊変時計とは別である。天水線を基準にO側へのずれを見ると熱水反応を検討できる。

\Needspace{4\baselineskip}\paragraph{選択肢の検討（査読待ち初稿）}\mbox{}\par

\noindent\textbf{1：不適切} 重水素2Hは半減期約12.3年。。2H（D）は安定同位体である。約12.3年でβ壊変するのは3H（T）であり、核種の質量数を取り違えている。\par

\noindent\textbf{2：適切（条件付き）} 一般に降水・陸水は海水より重い同位体に乏しい。。海からの蒸発と降水の履歴により一般に陸水は重同位体に乏しい。蒸発湖等の例外を分け、全陸水の厳密不等式ではない。\par

\noindent\textbf{3：適切（条件付き）} 反応で同位体組成の偏りが生じる作用は同位体分別。。同じ元素の同位体が反応・相変化で異なる比で配分されることである。核壊変で元素そのものが変わる作用とは区別する。\par

\noindent\textbf{4：適切（条件付き）} δ2H=\allowbreak{}8δ18O+10が天水線。。全球的な経験関係で、地域ごとの傾き切片・蒸発混合の逸脱がある。全水で厳密に成り立つ式ではない。\par

\noindent\textbf{5：適切（条件付き）} 熱水でδ2Hほぼ一定、δ18Oだけ高くなる場合。。高温水岩交換で岩石の豊富なOと水のOが交換し、H変化よりO変化が目立つ例がある。マグマ性流体混合など他機構も検討する。\par

\paragraph{再利用する考え方} 核種の生成、試料への取り込み、移動・分別、閉鎖、壊変を順番に分ける。比の基準と時計の開始条件を特定し、測定値と解釈する年代を区別する。

\paragraph{共通理論} \hyperref[theory:isotope:delta]{第\ref*{theory:isotope:delta}節}, \hyperref[theory:isotope:fractionation]{第\ref*{theory:isotope:fractionation}節}, \hyperref[theory:isotope:meteoric]{第\ref*{theory:isotope:meteoric}節}, \hyperref[theory:isotope:hydrothermal]{第\ref*{theory:isotope:hydrothermal}節}, \hyperref[theory:isotope:tritium]{第\ref*{theory:isotope:tritium}節}

\par 参考：\cite{iso_tritium2023,iso_vsmow2,gc_meteoric,iso_hydrothermal}。

\paragraph{関連問題} \hyperref[q:h27-reg-applied-25]{H27-REG-APPLIED-25}

\par 難易度 3（校準初稿）、優先度未評価。

\subsection{令和6年度 問21}
\label{q:r06-reg-applied-21}

\noindent\textbf{原典：} sources/R06/17\_applied\_science.pdf、PDF物理7頁。

\textbf{公式正答：} 1（single）。

\par\textbf{検討未確定：} 公式と初回は一致。原典の硫黄「国際標準物質CDTを用いて表記」は、出題前1995年のVCDT/IAEA-S-1勧告と物質・尺度の区別を省く。旧名称を現在の測定用標準物質へそのまま昇格しない。氷期δ大小の誤りは独立に説明できるが、硫黄肢の文字通りの妥当性は査読事項。

\paragraph{問題の要旨} 氷期と間氷期の有孔虫酸素比の大小を、蒸発・炭素混合・歴史的硫黄尺度から区別する。

\Needspace{4\baselineskip}\paragraph{確認した条件・記述（独自要約）}\mbox{}\par

\noindent\textbf{方向：} 最も不適切なもの。\par

\noindent\textbf{対象：} 同位体の標準・分別又は年代。番号・核種・旧標準名称を原典画像で確認。\par

\paragraph{問われる核心} 氷床・炭酸塩分別と同位体尺度の時点

\paragraph{必要知識} 海洋と氷の同位体収支、温度依存の炭酸塩分別、古い炭素混合、CDTとVCDT・IAEA-S-1の違い。

\paragraph{考え方・正答を選ぶ理由（査読待ち初稿）} 氷期より間氷期のδ18Oが大きいとする1が意図された不適切肢である。氷床へ軽い水を隔離する作用と、低温で殻が重い酸素を取り込みやすい分別が、通常ともに氷期の殻δ18Oを増す。他肢の条件を分け、硫黄CDTは歴史的起源として説明しVCDT尺度の実現を併記する。

\Needspace{4\baselineskip}\paragraph{選択肢の検討（査読待ち初稿）}\mbox{}\par

\noindent\textbf{1：意図不適切（大小逆）} 有孔虫δ18Oは氷期より間氷期に大きい。。氷期は16Oに富む水が陸上氷に隔離され、海水が18Oに富む。低温ほど炭酸塩と水の分別も殻18Oを増やす方向であり、有孔虫δ18Oは通常氷期に大きい。局所塩分・水温・生物効果・続成変質の補正は別に要する。原典はこの大小を逆にしている。\par

\noindent\textbf{2：適切（条件付き）} 有機物分解の炭素が地下水14C見かけ年代を古くする場合。。14Cに乏しい古い有機物又は炭酸塩から炭素が加わると、炭素量に対する14Cが薄まり未補正見かけ年代は古くなる。「ことがある」という原典の条件付き主張である。若い有機炭素なら若返る場合もあり、全有機物に一般化しない。\par

\noindent\textbf{3：適切（条件付き）} 蒸発残留水はDと18Oに富む。。軽い水分子が蒸発しやすく、残った水はD（2H）・18Oに富む傾向となる。供給水や湿度・同位体交換による例外を分ける。蒸気が重くなるという説明へ逆転しない。\par

\noindent\textbf{4：適切（条件付き）} 海水δ18Oの変動は天水・熱水より小さく後者は負値が多い。。外洋海水は天水に比べ同位体組成の変動が小さいという一般的比較である。蒸発・降水・氷融解・河川流入で地域差はある。熱水も天水起源なら負値を取り得るが、水岩交換・マグマ性流体で変わるため全海水・全熱水の固定値とはしない。\par

\noindent\textbf{5：旧尺度として意図適切・名称疑義} 硫黄標準はCDTトロイライト。。CDTはCanyon Diablo隕石のトロイライトFeSに由来する硫黄の歴史的基準。実際のVCDT実現は硫化銀IAEA-S-1等を使い、旧CDT物質を無条件に現行標準としない。\par

\paragraph{再利用する考え方} 核種の生成、試料への取り込み、移動・分別、閉鎖、壊変を順番に分ける。比の基準と時計の開始条件を特定し、測定値と解釈する年代を区別する。

\paragraph{共通理論} \hyperref[theory:isotope:climate]{第\ref*{theory:isotope:climate}節}, \hyperref[theory:isotope:fractionation]{第\ref*{theory:isotope:fractionation}節}, \hyperref[theory:isotope:deadcarbon]{第\ref*{theory:isotope:deadcarbon}節}, \hyperref[theory:isotope:standards]{第\ref*{theory:isotope:standards}節}

\par 参考：\cite{iso_oxygen,iso_deadcarbon,iso_standards1995,iso_vcdt,iso_hydrothermal}。

\paragraph{関連問題} \hyperref[q:r02-reg-applied-20]{R02-REG-APPLIED-20}

\par 難易度 3（校準初稿）、優先度未評価。
